% ------------------------------------------------------------------------------
% Abstract (English)
% ------------------------------------------------------------------------------
\begin{center}
\textbf{\large ABSTRACT}\\
\end{center}
\vspace{1cm}

\begingroup
\fontsize{11pt}{13pt}\selectfont
\singlespacing

\textit{Using large language models as automated evaluators (LLM-as-a-Judge) to compare response pairs reduces the time and cost of manual evaluation, yet model decisions remain vulnerable to answer order reversals (position bias) and instructional changes (prompt sensitivity). In Indonesian, reliability testing of LLM-as-a-Judge remains limited because prior evaluations only measured single-order accuracy against human preferences without testing answer position and prompt variations. This study aims to measure the decision consistency of LLM-as-a-Judge against human preferences on Indonesian responses and to analyze the effects of answer position and prompt variations on its decision consistency. This experimental study evaluates 519 unique response pairs across ten domains from the IndoPref dataset using the Qwen3-32B model in non-thinking mode. Each sample is tested across six condition combinations (3,114 inferences) consisting of two answer position orders and three Indonesian prompt variations using holistic, three-criteria, and anti-bias evaluation strategies. Each factor is analyzed separately using accuracy, Cohen's kappa, position consistency, preference flip rate, first-position preference rate, prompt consistency, prompt flip rate, 95\% confidence intervals, and the McNemar test to test the hypotheses that answer position and prompt variations alter model decisions and affect consistency against human preferences.}

\vspace{0.5cm}
\noindent \textbf{Keywords:} \textit{human preference, Indonesian, LLM-as-a-Judge, position bias, prompt sensitivity}
\endgroup